\origpage{107}{N. Robertson, P. Seymour / Journal of Combinatorial Theory, Series B 89 (2003) 77--108}
\noindent of a member of $\mathscr{C}$ by the class of all partial $\Omega$-patchworks of tree-width $\le 3\theta + 1$. However, the latter differs from $\mathscr{C}'$, because the members of $\mathscr{C}'$ are robust. To show that $Q$ ($= (G_0, \mu_0, \Delta_0, \phi_0)$ say) is an enlargement of a member of $\mathscr{C}$ by $\mathscr{C}'$, we must show that $\Delta_0(e)$ is a robust patch for every $e \in E(Q)$ which is not an edge of $P|\mathscr{L}$. Actually, we shall prove more, that $\Delta_0(e)$ is robust for every $e \in E(Q)$. Let $e \in E(Q)$. If $e \in E(P)$ then $\Delta_0(e) = \Delta(e)$ and hence is robust since $P$ is robust, as required. We assume then that $e \notin E(P)$. Since $Q = P|\mathscr{L}^*$, it follows that $G_0$ is obtained from $M(G, \mathscr{L}^*)$ by adding a new edge $e(A)$ for each $A \in \mathscr{L}^*$, where $e(A)$ has set of ends $\bar{\pi}(A)$; and $\Delta_0(e(A))$ is the set of all groupings feasible in $P|A$. Since $e \notin E(P)$ and hence $e \notin E(M(G, \mathscr{L}^*))$, it follows that $e = e(A)$ for some $A \in \mathscr{L}^*$. Let $B = G \setminus A$; then $(A^-, B^-) \in \mathscr{L}^{*-} = \mathscr{L}_t$. Hence $(A^-, B^-)$ is robust by (2)(b). Let $X_1, X_2 \subseteq \bar{\pi}(A) = V(\Delta_0(e))$ with $|X_1| = |X_2|$ say, and $X_1 \cap X_2 = \emptyset$. Define $k = |\bar{\pi}(A)| - |X_1|$. Let $(H_1, H_2)$ be any separation of $A^-$ such that $\bar{\pi}(A) \setminus X_i \subseteq V(H_i)$ for $i = 1, 2$. Since $(A^-, B^-)$ is robust, there exists $i \in \{1, 2\}$ such that $|V(H_i) \cap V(H_j \cup B^-)| \ge |V(A^- \cap B^-)|$, where $j = 3 - i$. Subtracting $|V(H_i \cap B) \setminus V(H_j)|$ from both sides gives
\begin{equation*}
|V(H_1 \cap H_2)| \ge |V(B^- \cap H_j)| \ge |\bar{\pi}(A) \setminus X_j| = k,
\end{equation*}
since $\bar{\pi}(A) \setminus X_j \subseteq V(B^- \cap H_j)$. From Theorem 6.1 of [2], applied to $P|A$, there is a realization of $P|A$ such that $k$ of its components have non-empty intersection with both $\bar{\pi}(A) \setminus X_1$ and $\bar{\pi}(A) \setminus X_2$. Therefore there is a pairing with vertex set $\bar{\pi}(A)$, feasible in $P|A$, which pairs $X_1, X_2$. Since $\Delta_0(e)$ is the set of all groupings feasible in $P|A$, it follows that $\Delta_0(e)$ is robust, as required. This proves (3).

Let $\mathscr{D}'$ be the class of all members $(G, \mu, \Delta, \phi) \in \mathscr{D}$ such that $G$ has a tangle of order $\ge \frac{4}{3}\theta$. We have shown then that

\hspace*{12pt}(4) \textit{For all $P = (G, \mu, \Delta, \phi) \in \mathscr{D}'$, there is a rotund, proper tree-decomposition $(T, \tau)$ of $G^-$ such that all pieces of $P$ under $(T, \tau)$ belong to $\mathscr{C}^*$.}

By (1), (4) and 6.6, $\mathscr{D}'$ is well-quasi-ordered by simulation. If $P = (G, \mu, \Delta, \phi) \in \mathscr{D} \setminus \mathscr{D}'$ then $G^-$ has tree-width $\le 2\theta$ by 2.2, and hence so does $P$ since $P$ is rootless. By 6.2, $\mathscr{D} \setminus \mathscr{D}'$ is well-quasi-ordered by simulation, and hence so is $\mathscr{D} = \mathscr{D}' \cup (\mathscr{D} \setminus \mathscr{D}')$. This proves 6.7. $\square$

\sectionhead{7. More on isolation}

Here is a useful way to prove that locations $\theta$-isolate tangles. Let $\lambda$ be a tie-breaker in a hypergraph $G$, let $\mathscr{T}$ be a tangle in $G$, and let $(A, B) \in \mathscr{T}$. We say that $(A, B)$ is $\lambda$-\textit{linked to} $\mathscr{T}$ if there is no $(A', B') \in \mathscr{T}$ with smaller $\lambda$-order with $A \subseteq A'$ and $B' \subseteq B$.

\textbf{7.1.} \textit{Let $\mathscr{T}$ be a tangle of order $\ge \theta \ge 1$ in a hypergraph $G$ with a tie-breaker $\lambda$, and let $\mathscr{L} \subseteq \mathscr{T}$ be a location of order $< \theta$, every member of which is $\lambda$-linked to $\mathscr{T}$. Then $\mathscr{L}$ $\theta$-isolates $\mathscr{T}$.}
